\documentclass{standalone}
\usepackage{circuitikz}
\usetikzlibrary{calc}
\definecolor{circuitnotemark}{HTML}{FE00FE}
\makeatletter
\pgfdeclareshape{smacoax}{
  \anchor{center}{\pgfpointorigin}
  \anchor{text}{\pgfpointorigin}
  \anchor{north}{\pgfpoint{0cm}{0.3cm}}
  \anchor{south}{\pgfpoint{0cm}{-0.3cm}}
  \anchor{east}{\pgfpoint{0.3cm}{0cm}}
  \anchor{west}{\pgfpoint{-0.3cm}{0cm}}
  \anchor{pin 1}{\pgfpoint{-0.7cm}{0cm}}
  \anchor{pin 2}{\pgfpoint{0cm}{-0.7cm}}
  \anchor{bpin 1}{\pgfpoint{-0.3cm}{0cm}}
  \anchor{bpin 2}{\pgfpoint{0cm}{-0.3cm}}
  \backgroundpath{
    \pgfpathmoveto{\pgfpoint{-0.274cm}{-0.122cm}}
    \pgfpatharc{204}{516}{0.3cm}
    \pgfpathmoveto{\pgfpoint{-0.7cm}{0cm}}\pgfpathlineto{\pgfpointorigin}
    \pgfpathmoveto{\pgfpoint{0cm}{-0.7cm}}\pgfpathlineto{\pgfpoint{0cm}{-0.3cm}}
  }
  \foregroundpath{
    \pgfpathcircle{\pgfpointorigin}{0.07cm}
    \pgfusepath{fill}
  }
}
\makeatother
\begin{document}
\begin{circuitikz}[american, line width=0.8pt]
\ctikzset{bipoles/length=1.2cm}
\ctikzset{grounds/scale=1.36}
\coordinate (b2) at (2,-2);
\coordinate (b4) at (6,-2);
\coordinate (b6) at (10,-2);
\coordinate (b8) at (14,-2);
\coordinate (b10) at (18,-2);
\coordinate (d6) at (10,-6);
\coordinate (d8) at (14,-6);
\coordinate (b12) at (22,-2);
\coordinate (c2) at (2,-4);
\coordinate (c12) at (22,-4);
\coordinate (b3) at (4,-2);
\coordinate (b11) at (20,-2);
\coordinate (d3) at (4,-6);
\coordinate (d11) at (20,-6);
\node[smacoax, draw, xscale=-1] (part-J1) at (b2) {}; % line 3
\node[anchor=south] at (part-J1.north) {$J_{1}$}; % line 3
\draw (b4) to[C, l_=$C_{1}$, a^=$\mathrm{20fF}$] (b6); % line 4
\draw (b6) to[L, l_=$L_{1}$, a^=$12.665\,\mathrm{m}\mathrm{H}$] (b8); % line 5
\draw (b8) to[R, l_=$R_{1}$, a^=$20\,\Omega$] (b10); % line 6
\draw (d6) to[C, l_=$C_{\mathrm{p}}$, a^=$5\,\mathrm{p}\mathrm{F}$] (d8); % line 7
\node[smacoax, draw] (part-J2) at (b12) {}; % line 8
\node[anchor=south] at (part-J2.north) {$J_{2}$}; % line 8
\node[ground] at (c2) {}; % line 9
\node[ground] at (c12) {}; % line 10
\draw (part-J1.pin 1) -- (b3); % line 12
\draw (b3) -- (b4); % line 12
\draw (b10) -- (b11); % line 13
\draw (b11) -- (part-J2.pin 1); % line 13
\draw (b3) |- (d6); % line 14
\draw (d8) -| (b11); % line 15
\draw (part-J1.pin 2) -- (c2); % line 16
\draw (part-J2.pin 2) -- (c12); % line 17
\node[circ] at (b3) {};
\node[circ] at (b11) {};
\path (1.619,-0.593) rectangle (2.381,0.593); % line 19
\node[anchor=west, inner sep=0, circuitnotemark, font=\large] at (2,0) {X}; % line 19
\path (21.619,-0.593) rectangle (22.381,0.593); % line 20
\node[anchor=west, inner sep=0, circuitnotemark, font=\large] at (22,0) {X}; % line 20
\node[anchor=south west, inner sep=2pt, circuitnotemark, font=\LARGE\bfseries] (circuittitle) at (current bounding box.north west) {X};
\path (circuittitle.south west) rectangle ++(8.138,0.853);
\end{circuitikz}
\end{document}
